\scanpage{553}
\mspage{464}

4.16. \uline{Exemples et questions}

4.16.\deleted{\unsure{1}}0. J'avais crû que tout foncteur ``raisonnable'' (\unsure{ou|et} plus précisément, \unsure{tout} foncteur ``logiquement accessible'') entre catégories accessibles devait être accessible. Ce n'est malheureusement pas vrai, puisqu'un composé de deux foncteurs Hom
\[
\mathrm{Ens}\longrightarrow\mathrm{Ens}\qquad
E\to\underbrace{\mathrm{Hom}(\mathrm{Hom}(E,I),J)}_{J^{(I^E)}}
\]
est non accessible, si $I,J$ sont tous deux de cardinal $\geq2$. Prenant \unsure{p.ex.|p. ex.} $I=J=\{0,1\}$, on trouve le foncteur
\[
E\longmapsto\unsure{\mathscr P|\mathcal P|\mathfrak P}(\mathscr P(E))
\]
qui n'est pas accessible.

4.16.1. Deuxième désillusion : j'avais crû que pour \unsure{tte|tt} catégorie accessible $\unsure{\mathcal M|M}$, \unsure{toute|tte} sous-catégorie \inserted{$\unsure{\mathcal N|N}$} \unsure{strict\textsuperscript{t}|strictem\textsuperscript{t}} pleine et stable par facteurs directs,
``raisonnable'', \deleted{\unsure{devait}} (p.ex. ``logiquement accessible'') devait être accessible dans $\unsure{M|\mathcal M}$. Il n'en est rien, \deleted{\unsure{comme on|comme}} même dans le cas particulier de l'image inverse de la sous-catégorie accessible \unsure{des|de} $\mathrm{Ens}$
% archive sheet number: 552
